% GENERATED FILE. Edit canonical lessons, then run npm run generate:books.
% canonical-source-hash: fnv1a64:7bb7a709f041aacd
% canonical-lessons: TA-C250-kutaram, TA-C250-uncal, TA-C250-alamaram, TA-C250-taracu, TA-C250-vaikkol

\chapter{Tent, Swing, Banyan tree, Pair of scales, Straw}
\label{ch:ta-tent-swing-banyan-tree-pair-of-scales-straw}
\hlchaptermodality{250}

\begin{chapteropening}
\textbf{By the end of this chapter:} \emph{I can say a tent, a swing, a banyan tree, a pair of scales and straw.}

Complete the last lesson of chapter 250: I can say a tent, a swing, a banyan tree, a pair of scales and straw.
\end{chapteropening}

\section[\texorpdfstring{\ta{கூடாரம்} (kūṭāram)}{கூடாரம் (kūṭāram)}]{\ta{கூடாரம்} (kūṭāram) — a tent}

\label{lesson:TA-C250-kutaram}



\begin{quote}
Before the new one: say the Tamil for chess, then the Tamil for a competition.
\end{quote}

\subsection*{You'll want to know: \ta{கூடாரம்}}
\textbf{\ta{கூடாரம்}} — \emph{kūṭāram} — \textquotedblleft{}a tent\textquotedblright{}.

\begin{grammarlens}[title={What you've built}]
One more word for this chapter.
\end{grammarlens}

\subsection*{Your turn}
\begin{itemize}
\raggedright
  \item \emph{Say it:} \emph{kūṭāram}
  \item \emph{Say it:} \emph{kūṭāram}, once more
  \item \emph{Say it:} say \emph{kūṭāram}
  \item \emph{Recall:} say the Tamil for chess, then the Tamil for a competition, then say \emph{kūṭāram} again
\end{itemize}

\begin{tcolorbox}[breakable,colback=teal!4,colframe=teal!35!black,title={Before you move on}]
What does \ta{கூடாரம்} mean? (\textquotedblleft{}A tent\textquotedblright{}.) Say it once more.
\end{tcolorbox}
\section[\texorpdfstring{\ta{ஊஞ்சல்} (ūñcal)}{ஊஞ்சல் (ūñcal)}]{\ta{ஊஞ்சல்} (ūñcal) — a swing}

\label{lesson:TA-C250-uncal}



\begin{quote}
Before the new one: say the Tamil for a competition, then the Tamil for a tent.
\end{quote}

\subsection*{You'll want to know: \ta{ஊஞ்சல்}}
\textbf{\ta{ஊஞ்சல்}} — \emph{ūñcal} — \textquotedblleft{}a swing\textquotedblright{}.

\begin{grammarlens}[title={What you've built}]
One more word for this chapter.
\end{grammarlens}

\subsection*{Your turn}
\begin{itemize}
\raggedright
  \item \emph{Say it:} \emph{ūñcal}
  \item \emph{Say it:} \emph{ūñcal}, once more
  \item \emph{Say it:} say \emph{ūñcal}
  \item \emph{Recall:} say the Tamil for a competition, then the Tamil for a tent, then say \emph{ūñcal} again
\end{itemize}

\begin{tcolorbox}[breakable,colback=teal!4,colframe=teal!35!black,title={Before you move on}]
What does \ta{ஊஞ்சல்} mean? (\textquotedblleft{}A swing\textquotedblright{}.) Say it once more.
\end{tcolorbox}
\section[\texorpdfstring{\ta{ஆலமரம்} (ālamaram)}{ஆலமரம் (ālamaram)}]{\ta{ஆலமரம்} (ālamaram) — a banyan tree}

\label{lesson:TA-C250-alamaram}



\begin{quote}
Before the new one: say the Tamil for a tent, then the Tamil for a swing.
\end{quote}

\subsection*{You'll want to know: \ta{ஆலமரம்}}
\textbf{\ta{ஆலமரம்}} — \emph{ālamaram} — \textquotedblleft{}a banyan tree\textquotedblright{}.

\begin{grammarlens}[title={What you've built}]
One more word for this chapter.
\end{grammarlens}

\subsection*{Your turn}
\begin{itemize}
\raggedright
  \item \emph{Say it:} \emph{ālamaram}
  \item \emph{Say it:} \emph{ālamaram}, once more
  \item \emph{Say it:} say \emph{ālamaram}
  \item \emph{Recall:} say the Tamil for a tent, then the Tamil for a swing, then say \emph{ālamaram} again
\end{itemize}

\begin{tcolorbox}[breakable,colback=teal!4,colframe=teal!35!black,title={Before you move on}]
What does \ta{ஆலமரம்} mean? (\textquotedblleft{}A banyan tree\textquotedblright{}.) Say it once more.
\end{tcolorbox}
\section[\texorpdfstring{\ta{தராசு} (tarācu)}{தராசு (tarācu)}]{\ta{தராசு} (tarācu) — a pair of scales, a balance}

\label{lesson:TA-C250-taracu}



\begin{quote}
Before the new one: say the Tamil for a swing, then the Tamil for a banyan tree.
\end{quote}

\subsection*{You'll want to know: \ta{தராசு}}
\textbf{\ta{தராசு}} — \emph{tarācu} — \textquotedblleft{}a pair of scales, a balance\textquotedblright{}.

\begin{grammarlens}[title={What you've built}]
One more word for this chapter.
\end{grammarlens}

\subsection*{Your turn}
\begin{itemize}
\raggedright
  \item \emph{Say it:} \emph{tarācu}
  \item \emph{Say it:} \emph{tarācu}, once more
  \item \emph{Say it:} say \emph{tarācu}
  \item \emph{Recall:} say the Tamil for a swing, then the Tamil for a banyan tree, then say \emph{tarācu} again
\end{itemize}

\begin{tcolorbox}[breakable,colback=teal!4,colframe=teal!35!black,title={Before you move on}]
What does \ta{தராசு} mean? (\textquotedblleft{}A pair of scales, a balance\textquotedblright{}.) Say it once more.
\end{tcolorbox}
\section[\texorpdfstring{\ta{வைக்கோல்} (vaikkōl)}{வைக்கோல் (vaikkōl)}]{\ta{வைக்கோல்} (vaikkōl) — straw, hay}

\label{lesson:TA-C250-vaikkol}



\begin{quote}
Before the new one: say the Tamil for a banyan tree, then the Tamil for a pair of scales.
\end{quote}

\subsection*{You'll want to know: \ta{வைக்கோல்}}
\textbf{\ta{வைக்கோல்}} — \emph{vaikkōl} — \textquotedblleft{}straw, hay\textquotedblright{}.

\begin{grammarlens}[title={What you've built}]
One more word for this chapter.
\end{grammarlens}

\subsection*{Your turn}
\begin{itemize}
\raggedright
  \item \emph{Say it:} \emph{vaikkōl}
  \item \emph{Say it:} \emph{vaikkōl}, once more
  \item \emph{Say it:} say \emph{vaikkōl}
  \item \emph{Recall:} say the Tamil for a banyan tree, then the Tamil for a pair of scales, then say \emph{vaikkōl} again
\end{itemize}

\begin{tcolorbox}[breakable,colback=teal!4,colframe=teal!35!black,title={Before you move on}]
What does \ta{வைக்கோல்} mean? (\textquotedblleft{}Straw, hay\textquotedblright{}.) Say it once more.
\end{tcolorbox}
